\chapter{Introdução}

Estudos recentes indicam que a aviação mundial entrou em um novo ciclo de crescimento sustentado, superando os níveis observados no período pré-pandêmico \cite{panrotas2026}. No contexto latino-americano, o Brasil tem desempenhado um papel de destaque, sendo responsável por uma parcela significativa do aumento no transporte aéreo de passageiros, com crescimento expressivo na demanda \cite{aeroin2026, portalin2026}.

Esse cenário é observado em diferentes regiões do país, com crescimentos relevantes no número de passageiros transportados, como no Nordeste e Sul \cite{mercadoeventos2026, r72026}. A manutenção de uma sequência prolongada de meses de alta no setor evidencia não apenas a recuperação após a pandemia, mas também uma tendência de expansão estrutural da aviação regional.

Além do crescimento observado no transporte de passageiros, o setor de carga aérea também tem apresentado uma expansão significativa nos últimos anos. Impulsionado por mudanças na logística global, aumento do comércio eletrônico e necessidade de cadeias de suprimento mais ágeis, o transporte aéreo de cargas vem registrando volumes recordes e projeções consistentes de crescimento \cite{agenciainfra2025, gazetadopovo2026}.

Esse avanço reforça a importância estratégica da aviação cargueira, especialmente em mercados regionais, onde a flexibilidade operacional e a capacidade de integração logística são fatores essenciais. Nesse contexto, aeronaves com capacidade mista ou adaptáveis a diferentes perfis de missão tornam-se particularmente atrativas, permitindo atender tanto à demanda por transporte de passageiros quanto ao escoamento eficiente de cargas em rotas de menor densidade.

Diante desse contexto, o presente trabalho tem como objetivo o desenvolvimento do projeto conceitual de uma aeronave regional, considerando tanto o transporte de passageiros, quanto de carga. O desenvolvimento inicial será fundamentado em estudos de mercado e em análises comparativas, permitindo a definição de um ponto de partida consistente. Essa abordagem possibilita a obtenção de estimativas preliminares para parâmetros fundamentais, como geometria externa e interna, configuração de cabine, arranjo estrutural e dimensões principais da aeronave. Além disso, serão avaliadas diferentes configurações aerodinâmicas, incluindo o tipo de asa, dispositivos hipersustentadores (flaps), superfícies de controle e empenagens, de modo a selecionar soluções que conciliem desempenho, eficiência e viabilidade operacional.

Paralelamente, o projeto será desenvolvido com base em metodologias consolidadas da literatura especializada, utilizando obras de referência reconhecidas na área de projeto aeronáutico. Esses materiais fornecem suporte teórico e prático para a realização das estimativas de desempenho, dimensionamento preliminar e análise de viabilidade da aeronave, além de reforçar o caráter evolutivo do projeto aeronáutico, no qual novas configurações derivam e aprimoram soluções previamente estabelecidas \cite{gudmundsson2021, raymer2018, anderson1999}.